\documentclass[11pt,a4paper]{article}
\usepackage[margin=22mm]{geometry}\usepackage[T1]{fontenc}
\usepackage{lmodern,amsmath,amssymb,amsthm,booktabs,microtype,fancyhdr}
\usepackage[colorlinks=true,linkcolor=blue,urlcolor=blue]{hyperref}
\hypersetup{pdftitle={Gauge-invariant certification of critical-orbit preparation},pdfauthor={Unified Dynamics research workstream}}
\setlength{\parindent}{0pt}\setlength{\parskip}{6pt}\setlength{\headheight}{14pt}
\pagestyle{fancy}\fancyhf{}\lhead{Unified Dynamics | Preparation certificate}\rhead{GPC v0.1}\cfoot{\thepage}
\newtheorem{theorem}{Theorem}\newtheorem{lemma}[theorem]{Lemma}\newtheorem{corollary}[theorem]{Corollary}
\newcommand{\tr}{\operatorname{tr}}\newcommand{\diag}{\operatorname{diag}}\newcommand{\Rea}{\operatorname{Re}}\newcommand{\Om}{\mathcal O_*}\newcommand{\F}{\mathrm F}
\begin{document}
\begin{center}
{\Large\bfseries Gauge-invariant certification\\of critical-orbit preparation}\\[9pt]
{\large A spanning-tree distance bound and a loop-only obstruction}\\[9pt]
Unified Dynamics research workstream\quad |\quad 18 October 2026 UTC\\
Version 0.1\quad |\quad Derived conditional theorem
\end{center}
\begin{abstract}
We give a global, constructive upper bound on distance to the frozen neutral critical orbit of the finite tetrahedral gauge model. The bound uses 27 real gauge-invariant scalars, including electric energy, full cycle traces and transported matter coherences. It vanishes exactly on that orbit and has an explicit local conditioning bound. Rational interval inputs can therefore certify entry into the earlier nonlinear trapping ball without assuming beforehand that the state is in its local chart. Conversely, active matter-dressed loops cannot certify full preparation: zero-Gauss electric perturbations are invisible instantaneously at every rank, and unused-color dynamics remain invisible for all time at rank at least two. The measurement contract is mathematical and conditional; no experimental realization or physical correspondence is claimed.
\end{abstract}
Write $d=d_g(x,\Om)$ in the frozen unnormalized Frobenius/unit-sphere metric. Section 1 defines $s,u\in[0,1]$, fourteen residuals $h_e$ and $\mathcal E=\sum_e\|E_e\|_\F^2$.
\begin{theorem}[Constructive preparation certificate]\label{th:main}
At each finite rank $r\ge1$, for every fixed-norm state with $s>0$, define
\[
C=\sum_{e\notin\mathcal T}h_e^2+2(1-s)+2(1-u),\qquad
\boxed{\mathcal P=\mathcal E+\frac{121}{49}C.}
\]
Then $d^2\le\mathcal P$ globally, and $\mathcal P=0$ if and only if $x\in\Om$. If $d\le1/8$, then $s\ge3/4$ and
\[
\boxed{d^2\le\mathcal P\le\frac{442425}{49}d^2\le10000d^2.}
\]
All quantities are invariant under node gauge transformations and independent species phases. The statements hold on the full fixed-norm manifold and hence on its zero-Gauss subset.
\end{theorem}
\textbf{Status.} Development lane, no registry entry; constructor verification only; independent review pending; physical promotion $0$. The reference action, metric, coefficients and ledger are unchanged. This certificate separates one frozen orbit; it is not complete state tomography.
\newpage
\section{Frozen carrier, orbit and the spanning tree}
The carrier has the fifteen nonempty faces of a tetrahedron, ordered by increasing size and then lexicographically on vertices $0,1,2,3$. There are 28 links $e=(i\leftarrow j)$ with $i$ a codimension-one face of $j$. The boundary sign is $b_e=(-1)^h$ when the $h$th ordered vertex is removed. There are eighteen diamonds. The critical representative is
\[
U_e^*=\diag(-ib_e,I_{r-1}),\qquad E_e^*=0,\qquad q_i^*=w_i^*=v_i e_1,
\]
where $v>0$ is the unit Perron vector of unsigned Hasse adjacency, with eigenvalue $\rho=\sqrt{8+2\sqrt{10}}$. The global matter vectors have norm one. The orbit $\Om$ is generated by $U(r)^{15}$ and two independent species phases. Its distance uses the product of bi-invariant Frobenius link metrics, Euclidean Hermitian electric metrics and real round matter-sphere metrics [1,2].

The frozen Hamiltonian, with $a,N>0$, $\kappa=rNk_*$, is
\[
H=\frac a2\sum_e\tr E_e^2+\kappa\sum_d(1-\Rea\tr W_d/r)
+N(q^\dagger h_Uq+w^\dagger h_Uw)+2\mu N.
\]
Here $D_U[i,j]=b_eU_e$, $h_U=i(D_U^\dagger-D_U)$; the two species have opposite symplectic forms, and $W_d$ is defined in Section 6. Electric matrices are Hermitian.

Root the tree $\mathcal T$ at the tetrahedral cell (index 14). For each proper face $i$, its parent $p(i)$ is obtained by adding its smallest missing vertex. This gives fourteen tree links and maximum root-to-node depth three. Define root-to-node transports and nominal active phases by
\[
T_{14}=I,\quad T_i=U_{i,p(i)}T_{p(i)},\qquad
t_{14}=1,\quad t_i=(-ib_{i,p(i)})t_{p(i)},\quad \eta_i=\bar t_i.
\]
The root-framed fields are $\widetilde q_i=T_i^\dagger q_i$, $\widetilde w_i=T_i^\dagger w_i$. For a non-tree edge put
\[
H_e=T_i^\dagger U_eT_j,\qquad
\zeta_e=\bar t_i(-ib_e)t_j\in\{-1,+1\}.
\]
Exactly ten of the fourteen signs are negative and four positive. Tree-framed nominal links are identity on tree edges and $I+(\zeta_e-1)P_0$ on other edges, with $P_0=e_1e_1^\dagger$.

Define the transported matter summaries and the root active projector:
\[
S=\sum_i v_i t_i\widetilde q_i,\qquad Z=\sum_i v_i t_i\widetilde w_i,
\qquad s=\|S\|>0,\quad n=S/s,\quad P=nn^\dagger,
\]
\[
u=|n^\dagger Z|,\qquad D_e(P)=I+(\zeta_e-1)P,\qquad
h_e=\|H_e-D_e(P)\|_\F.
\]
Cauchy--Schwarz gives $s,u\le1$. Under node gauge $g_i$, $T_i$ transforms as $g_iT_i g_{14}^\dagger$; hence $H_e$ transforms by root conjugation. The vectors $S,Z$ transform by the same root unitary, with their own species phases. Thus $P$ transforms by root conjugation and $s,u,h_e,\mathcal E$ are invariant. No choice of external root color frame is part of the certificate.
\newpage
\section{Global distance certificate: proof}
Apply the actual node gauge $T_i^\dagger$ to the state. Tree links become identity. A corresponding representative of $\Om$ has links $I$ on the tree, $D_e(P)$ elsewhere, zero electric field, and matter profiles
\[
f_i=v_i\eta_i n,\qquad e^{i\beta}f_i.
\]
To see membership in $\Om$, first tree-frame the nominal configuration using
$T_i^*=\diag(t_i,I_{r-1})$, then apply a common root unitary taking $e_1$ to $n$. Both operations are allowed node gauge transformations. The phase $\beta$ is an allowed phase of the second species.

The first species' real inner product with $f$ is
\[
\Rea\sum_i f_i^\dagger\widetilde q_i=\Rea(n^\dagger S)=s.
\]
The phase of the second profile may be chosen so that its real inner product with $\widetilde w$ is $|n^\dagger Z|=u$, including the case $u=0$. Therefore the squared matter chord distances to this representative are exactly $2(1-s)$ and $2(1-u)$.

For real unit vectors at sphere angle $\theta\in[0,\pi]$, their chord is $2\sin(\theta/2)$ and
\[
\theta\le\frac\pi2\,2\sin(\theta/2).
\]
The same constant applies to unitary matrices: diagonalize $U^\dagger V$, take principal eigenangles in $[-\pi,\pi]$ and sum the corresponding squared scalar inequalities. Thus unitary geodesic distance is at most $\pi/2$ times Frobenius chord distance. Since electric distance to zero is exactly $\sqrt{\mathcal E}$, the constructed orbit representative proves
\[
d^2\le\mathcal E+\frac{\pi^2}{4}
\left[\sum h_e^2+2(1-s)+2(1-u)\right].
\]
A self-contained rational majorant follows from
\[
0<\int_0^1\frac{x^4(1-x)^4}{1+x^2}\,dx=\frac{22}{7}-\pi,
\qquad \frac{\pi^2}{4}<\frac{121}{49}.
\]
The integral is strictly positive, and the equality follows by polynomial division and $\int_0^1(1+x^2)^{-1}dx=\pi/4$. This proves the global upper bound in Theorem~\ref{th:main}, without a local chart or a Gauss constraint.

Every term of $\mathcal P$ is nonnegative. If it vanishes, both matter chord distances, every non-tree link chord and every electric field vanish; tree links already agree. The state is the constructed orbit representative. Conversely, for any member of $\Om$, tree transports give exactly these common active profiles and links, so $s=u=1$, $h_e=0$, $\mathcal E=0$. This proves the zero-set assertion. If $s=0$, this particular construction is undefined and returns no certificate; it does not return evidence of instability.
\newpage
\section{Explicit local conditioning: proof}
The compact orbit has a closest point. By a common gauge/species-phase isometry, align this point with the stated representative. Let $L$ be the product link geodesic distance to that point, $Q,W$ the two matter chord distances, and $\mathcal E$ the squared electric distance. Then
\[
L^2+Q^2+W^2+\mathcal E\le d^2.
\]
Each link operator-norm difference is at most its geodesic distance. Telescoping a tree transport with at most three factors, using Cauchy--Schwarz and $\sum_i v_i^2=1$, gives
\[
\|\widetilde q-\widetilde q^*\|\le Q+\sqrt3L\le2d,
\quad \|\widetilde w-\widetilde w^*\|\le W+\sqrt3L\le2d.
\]
For example, the squared weighted transport error is bounded by
$3\sum_i v_i^2\sum_{e\in\operatorname{path}(i)}\ell_e^2\le3L^2$; no rank factor is introduced. Consequently $\|S-e_1\|\le2d$ and
\[
s\ge1-2d\ge3/4\quad(d\le1/8),\qquad
\|n-e_1\|\le2\|S-e_1\|\le4d.
\]
The normalization bound follows by adding $\|S-n\|=|s-1|$ to $\|S-e_1\|$. For unit vectors,
\[
\|nn^\dagger-e_1e_1^\dagger\|_\F
\le\sqrt2\|n-e_1\|\le6d.
\]
Each non-tree $H_e$ is a word with at most seven link factors, and any individual tree edge occurs at most twice. Its squared multiplicity sum is at most fourteen. Telescoping with unitarily invariant Frobenius norm therefore yields
\[
\|H_e-H_e^*\|_\F\le\sqrt{14}L\le4d,
\qquad \|D_e(P)-D_e(P_0)\|_\F\le2\|P-P_0\|_\F\le12d.
\]
Thus $h_e\le16d$. Each chosen matter-profile chord is at most $2d+4d=6d$; choosing the optimal second-species phase only improves this bound. We obtain
\[
C\le[14(16)^2+2(6)^2]d^2=3656d^2,
\]
\[
\mathcal P\le\left(1+\frac{121}{49}3656\right)d^2
=\frac{442425}{49}d^2\le10000d^2.
\]
This completes Theorem~\ref{th:main}. In particular, if an exact score satisfies $\mathcal P\le1/64$, its global upper bound first implies $d\le1/8$ and hence
\[
\sqrt{\mathcal P}/100\le d\le\sqrt{\mathcal P}.
\]
For an interval upper score $\widehat{\mathcal P}$, only the upper distance implication is automatic; its excess measurement uncertainty cannot be used as a lower distance bound.
\newpage
\section{Twenty-seven scalar inputs and guaranteed enclosures}
The construction can be evaluated from real gauge-invariant scalars
\[
\alpha=S^\dagger S=s^2,\quad \gamma=|S^\dagger Z|^2,
\quad\mathcal E,\quad t_e=\Rea\tr H_e,
\quad m_e=\Rea(S^\dagger H_eS)\quad(\zeta_e=-1).
\]
There are three global scalars, fourteen traces and ten active moments: 27 in total, at every fixed finite rank. This is an observable-count statement, not a claim about rank-independent acquisition cost. We have $0<\alpha\le1$, $0\le\gamma\le\alpha$ and
\[
s=\sqrt\alpha,\qquad u=\sqrt{\gamma/\alpha},\qquad
h_e^2=\begin{cases}
2r-2t_e,&\zeta_e=+1,\\
2r-2t_e+4m_e/\alpha,&\zeta_e=-1.
\end{cases}
\]
The trace identity follows by expanding $\|H_e-D_e(P)\|_\F^2$ and using $P=SS^\dagger/\alpha$ and unitary $H_e,D_e(P)$.

Suppose guaranteed rational enclosures provide
\[
0<a_L\le\alpha\le a_U\le1,\qquad
0\le g_L\le\gamma\le g_U\le1,\qquad\mathcal E\le e_U,
\]
with $e_U\ge0$, and ordered intervals $[t_{e,L},t_{e,U}]$, $[m_{e,L},m_{e,U}]$. Let $s_L\le\sqrt{a_L}$ and $u_L\le\sqrt{g_L/a_U}$ be nonnegative certified lower bounds. For a negative cycle set
\[
M_e=\max\{m/a:\ m\in\{m_{e,L},m_{e,U}\},\ a\in\{a_L,a_U\}\}.
\]
This four-corner maximum remains valid for either sign of $m$. Define
\[
b_e=\begin{cases}2r-2t_{e,L},&\zeta_e=+1,\\
2r-2t_{e,L}+4M_e,&\zeta_e=-1,\end{cases}
\]
\[
\boxed{\widehat{\mathcal P}=e_U+\frac{121}{49}
\left[\sum b_e+2(1-s_L)+2(1-u_L)\right]\ge\mathcal P\ge d^2.}
\]
The enclosure calculator uses exact fractions and integer square roots on a $10^{80}$ grid. It rejects nonpositive coherence lower bounds, reversed endpoints, negative electric bounds and negative derived residual upper bounds. Rejection is not a dynamical verdict. Passing the arithmetic check does not prove that the intervals come from a realizable state or actually cover the state: acquisition, calibration, admissibility and coverage are external inputs to this conditional theorem. Ordinary double precision subtraction near the identity cannot supply the required tiny interval widths.
\newpage
\section{Connection to the nonlinear trapping theorem}
At the frozen critical endpoint, $\kappa=rNk_*$, $a,N>0$, [1,2] provide a certified radius $R_{\rm p}$ and, for any $0<\eta\le R_{\rm p}$,
\[
c=N/1600,\qquad J=\max(a,25N),\qquad B=c\eta^4,
\qquad \delta^2=c\eta^4/J.
\]
Inside the tube, $H-H_*\ge c d^4$, and $H-H_*\le Jd^2/2$. The endpoint robustness theorem [1] admits invariant Hamiltonian errors split into generic and orbit-centered parts, with full-tube bounds $\|\nabla V_g\|\le\nu$, $\nabla V_c|_{\Om}=0$, $\|\operatorname{Hess}V_c\|\le\beta$ and
\[
\nu+\eta\beta\le c\eta^3/8.
\]
Its time-dependent extension also requires cumulative gradient variation at most $c\eta^3/8$ through the specified horizon.
\begin{corollary}[Observable sufficient entry condition]
Under all frozen action, norm, zero-Gauss and error hypotheses of [1], guaranteed scalar enclosures satisfying
\[
\widehat{\mathcal P}(x_0)\le\delta^2
\]
certify that the initial state belongs to its trapping ball. Consequently $d_g(x(t),\Om)<\eta$ for the time interval covered by [1]. No local state-chart assumption is needed to infer the initial distance.
\end{corollary}
\begin{proof}
The global certificate gives $d_g(x_0,\Om)\le\sqrt{\widehat{\mathcal P}}\le\delta$. Apply [1]. The preparation scalars do not certify its coefficient, Gauss or time-variation assumptions; those must be checked separately.
\end{proof}
A simple sufficient raw-interval allocation can be stated relative to the ideal record. Let $0\le\varepsilon<1$, $e_U\le\varepsilon$, $\alpha,\gamma\in[1-\varepsilon,1]$, every trace enclosure be centered at its nominal value with half-width at most $\varepsilon$, and every negative-cycle moment lie in $[-1,-1+\varepsilon]$. Nominal traces are $r$ on positive cycles and $r-2$ on negative cycles. Choose the exact rational lower bounds $s_L=u_L=1-\varepsilon$. Then
\[
b_e\le2\varepsilon\ (\zeta_e=+1),\quad
b_e\le6\varepsilon\ (\zeta_e=-1),\quad
\widehat{\mathcal P}\le\left(1+\frac{121}{49}(4\cdot2+10\cdot6+4)\right)\varepsilon
=\frac{8761}{49}\varepsilon\le179\varepsilon.
\]
Hence $\varepsilon\le\delta^2/180$ suffices with these analytic lower bounds. At $a=N=1$, the inherited inward radius is $\eta=1/829538022$ and $\delta^2\simeq5.27952465\times10^{-41}$, giving $\delta^2/180\simeq2.93306925\times10^{-43}$. This extremely conservative allocation is a sufficient model-unit precision budget, not a necessary precision, an experimental feasibility result or a conversion to physical units. The general integer-square-root implementation retains its rounding margin and tests its actual returned score.
\newpage
\section{Exact obstructions to loop-only preparation claims}
For a diamond $d=(i,j,h,k)$, let an integer edge circulation $z$ be $+1$ on $i\leftarrow j\leftarrow k$, $-1$ on $i\leftarrow h\leftarrow k$, zero elsewhere. Its node divergence is zero and $\sum_e z_e^2=4$. Active matter-dressed readouts are
\[
\mathfrak z_d=-q_k^\dagger W_dq_k/\|q_k\|^2,\qquad
W_d=(U_{ih}U_{hk})^\dagger U_{ij}U_{jk}.
\]
\begin{theorem}[Instantaneous configuration-only obstruction]
For every finite rank and any real $\lambda$, keep $U,q,w$ at their frozen values and set $E_e=\lambda z_e I_r$. This state has zero Gauss, unchanged configuration observables and $\mathfrak z_d=1$ on every diamond, while
\[
d_g(x,\Om)=2\sqrt r\,|\lambda|.
\]
Consequently no configuration-only observation can yield a finite global upper preparation-distance bound at this configuration without an additional electric, energy or dynamical assumption.
\end{theorem}
\begin{proof}
The fixed matter charges cancel between the two species. Scalar electric matrices commute with every link, so their Gauss contribution is exactly the zero node divergence of $z$. Every member of $\Om$ has $E=0$; the remaining coordinates already lie on the orbit. Electric product distance is $\sqrt{r\lambda^2\sum z_e^2}$ and is attained without moving other coordinates. Configuration observables do not involve $E$.
\end{proof}
\begin{theorem}[All-time unused-color obstruction]
For $r\ge2$ use $E_e=\lambda z_e e_2e_2^\dagger$ at the same configuration. Under the unchanged frozen Hamiltonian, the full history of active matter-dressed loops and active matter densities is identical to the nominal history, although the initial orbit distance is $2|\lambda|$ and is unbounded as $|\lambda|\to\infty$.
\end{theorem}
\begin{proof}
The initial Gauss argument still applies, since all links preserve the active/unused split. In the block diagonal sector with both species in the first color, the action is the sum of an active $U(1)$ matter/gauge Hamiltonian and an unused $U(r-1)$ pure-gauge Hamiltonian. The electric trace, Wilson trace and matter hopping terms split additively; their variations have no off-block source. The block diagonal subspace and active matter subspace are invariant. The active block starts on the frozen critical orbit, with its allowed species phase motion. The unused block evolves independently. Every active readout remains nominal for all time. Smooth finite-rank flow is complete here: compact link/matter variables and positive electric kinetic energy bound electric norms on each fixed finite-energy trajectory. Distance at initial time follows from $\sum z_e^2=4$.
\end{proof}
This is a no-go for the specified active-only observation map, not a no-go for all dynamic observers or an inference about the unrelated record-only sketch kernels. Full cycle traces and electric energy in the proposed certificate address precisely these hidden directions.
\newpage
\section{Requirements, verification and remaining work}
The input contract for a future experiment or emulator has five distinct obligations:
\begin{enumerate}
\item Establish the frozen carrier, finite gauge representation, normalized matter coordinates, metric convention and critical reference coefficients. Bound model discrepancy on the full trapping tube as required by [1].
\item Establish norm and Gauss admissibility separately, or derive a theorem that quantitatively admits their residual errors. The 27 scalars alone do not establish these constraints.
\item Obtain guaranteed joint-coverage enclosures for the transported coherences, full cycle traces, active cycle moments and electric squared norm. Assign a stated confidence or deterministic coverage interpretation; the present theorem assumes coverage rather than estimating it.
\item Evaluate a rigorous upper score, check it against $\delta^2$ and preserve its uncertainty arithmetic. Negative residual upper bounds or a vanishing coherence lower bound require correction of the input contract or a different certificate.
\item Supply the coefficient/drift bounds for the requested duration and an independently calibrated mapping to physical observables and units before interpreting this as a physical experiment requirement.
\end{enumerate}
This session's constructor replay supplies 46 exact checks: carrier/tree/sign counts, all eighteen circulation witnesses, word multiplicities, rational conditioning constants, the positive-integral rational majorant, inward square roots, accepted/rejected interval controls and the inherited unit budget. It retains unchanged copies of the endpoint and parity budget helpers and a hash-bound frozen endpoint certificate. These dependencies are conditional inputs, not newly established independent results.

Numerical diagnostics use eight full-matrix perturbation fixtures at ranks $1,2,3,5$ and scales $10^{-3},10^{-9}$, checking constructed-orbit geodesic bounds, the conservative path-length upper estimate and gauge/species-phase invariance. Four explicit hidden-electric controls are zero Gauss; three unused-color configuration controls preserve active loops while yielding a positive full score. Random perturbation fixtures need not be zero Gauss because the geometric upper theorem holds on the full manifold. Direct residual norms avoid catastrophic trace subtraction. Floating diagnostics do not prove the analytic inequalities or certify tiny acquisition intervals.

The next substantive work is an independent proof review and a less conservative preparation/energy certificate, followed by a scoped physical bridge and data-availability audit. No published experimental dataset has been reanalyzed in this milestone. No coupling unification, proton decay, monopole detection, continuum limit or empirical GUT result follows from the present finite model theorem.

\begin{thebibliography}{9}
\bibitem{ert} Unified Dynamics workstream. \emph{Robust trapping at the quartic critical endpoint}, ERT v0.1, 16 October 2026. Frozen capture: endpoint-robustness, UTC 2026-10-16T164203Z, suffix 80c0789d. Conditional inherited trapping and model-error theorem; review pending.
\bibitem{pet} Unified Dynamics workstream. \emph{Parity-sharpened endpoint trapping}, PET v0.1, 16 October 2026. Frozen capture: endpoint-parity, UTC 2026-10-16T161824Z, suffix 37748fe1. Conditional quartic tube and radius; review pending.
\bibitem{gov} UD Governance CRL2 Subdivision v1.0; Coefficient Provenance Ledger v0.39. Attached frozen governance and ledger; development lane, physical promotion zero.
\end{thebibliography}
\end{document}
